% Source: MQ14a_v1.html
\documentclass[11pt]{article}
\usepackage[margin=1in]{geometry}
\usepackage{amsmath}
\usepackage{amssymb}
\usepackage{siunitx}
\usepackage{enumitem}

\title{MQ14a: Oscillations}
\author{}
\date{}

\begin{document}
\maketitle

\section*{MQ14a: Oscillations}

\begin{enumerate}[leftmargin=*,label=\textbf{Question \arabic*.}]
\item \hfill \textit{1 pts}

Suppose a horizontal block-spring system has a spring constant 818 N/m and block of mass 12 kg. Calculate the angular frequency in rad/s of its harmonic motion..

\item \hfill \textit{1 pts}

A tunnel is drilled straight through the center of an airless planet to the other side. For simplicity, the planet is assumed to have constant density. The equation for the motion of an object falling through the tunnel is given by

$\frac{d^2r}{dt^2}+\left(\frac{4\pi G\rho}{3}\right)r=0$


where \emph{G} = 6.67 × 10$^{-11}$ m$^{3}$/(kgᐧ݀s$^{2}$) is the gravitation constant. The planet's density and radius are 𝜌 = 3,856 kg/m$^{3}$ and \emph{R} = 1,701 km.

If an astronaut jumps in the tunnel, how many \textbf{\emph{hours}} would it take for her to fall all the way to the other side and return to her original position?



\emph{T} = \rule{2cm}{0.4pt} h (Note the units!)

\item \hfill \textit{1 pts}

Fill In the Blanks with numerical values. \emph{\textbf{Give three digits for each answer}}, rounding as appropriate. For example, if your answer is "2", report "2.00". If your answer is "1.5555", report "1.56". If your answer is "0.7", report "0.700".



Suppose a mass-spring system has the description \emph{x}(\emph{t}) = 4.00 cos(2.00\emph{t} - 0.500), where each number has unshown MKS units.



The amplitude is
 m, the angular frequency is
 rad/s, the phase shift is
 rad, the frequency is
 Hz; the period is
 s.

\item \hfill \textit{1 pts}

Suppose a mass-spring system has spring constant 711 N/m and mass 6 kg. The mass is displaced in the negative x-direction by 0.43 m and set in motion at 2.59 m/s in the positive x-direction. Calculate the phase shift as a positive angle in radians, assuming a cosine solution.

\end{enumerate}

\end{document}
